\session{12}{11/14}{1限}{4}{デジタル戦略とAI戦略}

\goals{
  \item DX の意味を、単なるデジタル化と区別して説明できる。
  \item AIによるビジネスモデルの変化とプラットフォームの考え方を理解する。
  \item 自社開発と外部サービス利用の比較、AI投資の評価の考え方を説明できる。
}

\guidesec{解説}

\topic{DX とは何か}
DX（デジタルトランスフォーメーション）は、デジタル技術を使ってビジネスモデルや組織のあり方を変革し、競争優位を築くことです。
「紙をPDFにする」「会議をオンラインにする」といった作業のデジタル化は DX の出発点ではありますが、DX そのものではありません。
段階で整理すると次のようになります。
\par\noindent\begin{gtabh}{colspec={Q[l,wd=34mm] X[l] X[l]}}
  段階 & 内容 & 例 \\
  デジタイゼーション & アナログな情報や作業をデジタルに置き換える & 紙の伝票を電子化、FAX をメールに \\
  デジタライゼーション & 業務プロセス全体をデジタル技術で再設計する & 受注から請求までをシステムで一貫処理、第5回の自動化 \\
  デジタルトランスフォーメーション & 事業そのもの、顧客への価値の提供のしかた、組織を変える & 製品の販売からサブスクリプションへ、データを使った新サービス \\
\end{gtabh}
生成AIは、デジタライゼーションを加速させると同時に、顧客への価値の提供のしかた（個別化、対話型のサービス）を変えるため、DX の段階にも関わります。

\topic{AIによるビジネスモデルの変化}
ビジネスモデルとは、誰に、どんな価値を、どう届けて、どう収益を得るかの設計です。AIはこの各要素を変えます。
\begin{itemize}
  \item \textbf{パーソナライズ}　利用データから個々の顧客の好みや状況を推定し、一人ひとりに合わせた提案・価格・内容を提供する。「全員に同じものを売る」から「一人ひとりに合わせる」へ。
  \item \textbf{サービス化とサブスクリプション}　製品を売り切るのではなく、利用データを使って継続的にサービスを提供し、利用に応じて収益を得る。
  \item \textbf{対話型のインターフェース}　検索や画面操作の代わりに、自然言語での対話でサービスを提供する。顧客接点のあり方が変わる。
  \item \textbf{プラットフォーム}　売り手と買い手、作り手と使い手など複数の利用者グループを結びつけて価値を生む。参加者が増えるほど価値が高まるネットワーク効果が働き、先行者が有利になりやすい。AIは需要と供給のマッチングの精度を上げ、プラットフォームの価値を高める。
\end{itemize}
一方で、AIは参入障壁を下げる面もあります。誰でも一定水準の文章・画像・コードを作れるようになると、それらを強みにしていた事業の優位性は薄れます。
自社の強みがAIによって「誰でもできること」になっていないかを問い直す必要があります。

\topic{自社開発か、外部サービス利用か}
AIを事業に使うとき、自社で開発するか、外部のサービスを使うかは主要な選択です。実際には両者の間に段階があります。
\par\noindent\begin{gtabh}{colspec={Q[l,wd=32mm] X[l] X[l]}}
  選択肢 & 利点 & 課題 \\
  外部の SaaS をそのまま使う & 早く、安く始められる。提供側が更新する。 & 他社も使えるため差別化しにくい。データの扱いと機能に制約。 \\
  外部の AI モデルを API で使い、自社の仕組みに組み込む & 自社のデータや業務に合わせられる。 & 開発と運用の人材が必要。モデルの提供側への依存。 \\
  自社でモデルを調整・開発する & 差別化と、データの管理を自社で完結できる。 & 費用と人材の負担が大きい。技術の変化が速く、陳腐化のリスク。 \\
\end{gtabh}
判断の軸は、その業務が自社の競争優位の中核か（中核なら自社で、そうでなければ外部で）、自社固有のデータがあるか、
必要な人材と予算があるか、そしてデータの扱いに関する制約（第13回）です。多くの企業にとって現実的なのは、
外部サービスを基本にしつつ、競争優位に直結する部分だけを自社で作る組み合わせです。

\topic{AI投資の評価}
AI投資は、他の投資と同じく、コストと効果とリスクで評価します。ただし、次の特徴があります。
\begin{itemize}
  \item \textbf{コスト}　初期費用だけでなく、運用費（利用料、データ整備、人の確認の工数、教育）が継続的にかかる。
  \item \textbf{効果}　時間の削減のような測りやすい効果と、意思決定の質の向上や新サービスのような測りにくい効果がある。KPI を導入前に決める（第7回）。
  \item \textbf{リスク}　データの確保、法規制の変化、技術の陳腐化、誤りによる信用の失墜、提供側への依存。
  \item \textbf{段階的な投資}　不確実性が高いため、小さく始めて学びながら投資を広げる。PoC で判断基準を決めておく。
  \item \textbf{投資しないリスク}　競合が先行すること、人材が集まらなくなること、シャドーAIが広がることも、評価に含める。
\end{itemize}
「流行しているから」「競合が導入したから」だけで投資を決めると、目的が曖昧な技術先行の導入（第7回の失敗要因）になります。

\topic{キーワード}
\par\noindent\begin{keywords}{}
  DX & デジタル技術でビジネスモデルや組織を変革し、競争優位を築くこと。単なるデジタル化と区別する。 \\
  ビジネスモデル & 誰に、どんな価値を、どう届けて、どう収益を得るかの設計。 \\
  パーソナライズ & 利用データから個々の顧客に合わせた提案を行うこと。 \\
  プラットフォーム & 複数の利用者グループを結びつけて価値を生む事業。ネットワーク効果が働く。 \\
  自社開発／外部サービス & AIを使う際の選択。競争優位の中核か、固有データがあるか、人材があるかで決める。 \\
  AI投資の評価 & コスト・効果・リスクに加え、段階的な投資と「投資しないリスク」を考慮する。 \\
\end{keywords}

\guidesec{演習課題}

\exercise{1}{AI演習：選んだ企業のAI戦略を、AIを使って批判的に評価する}
\instr{AI活用を公表している企業を1社選び、統合報告書や中期経営計画、プレスリリースなど公開資料からAI戦略に関する部分を抜き出して、次のプロンプトで批判的に評価させてください。AIの指摘のうち、資料で確かめられるものと、AIの推測にすぎないものを分け、自分の評価を書いてください。}
\begin{promptbox}[AI戦略の批判的評価]
あなたは厳しい投資家です。次の資料は「（企業名）」のAI戦略に関する公開情報です。(1) この戦略が解決しようとしている経営課題は明確か、(2) 自社開発と外部サービスの使い分けは妥当か、(3) 効果をどう測るとしているか、(4) データ・人材・法規制のリスクにどう備えているか、(5) 競合と比べた差別化はあるか、の5点について、資料に根拠がある指摘と、資料にないため推測になる指摘を分けて評価してください。
（資料の抜粋）
\end{promptbox}

\exercise{2}{自社開発か外部サービスか}
\instr{次の3つの業務について、外部 SaaS・API 利用・自社開発のどれが適切かを、「競争優位の中核か」「固有データがあるか」「人材・予算」「データの制約」の4つの軸で判定し、理由を書いてください。企業は従業員500人の専門商社とします。}
\begin{itemize}
  \item (A) 社内の経費精算の問い合わせに答えるチャットボット
  \item (B) 30年分の取引データを使った、顧客ごとの最適な提案商品の推定
  \item (C) 取引先との契約書のリスク条項のチェック
\end{itemize}

\exercise{3}{投資評価のメモ}
\instr{演習2の (B) について、経営会議に出す投資評価のメモを「目的と KPI」「コスト（初期・運用）」「期待効果」「リスクと対策」「段階的な進め方」「投資しない場合のリスク」の6項目で書いてください。金額は仮の数値でかまいませんが、根拠（何を想定したか）を添えてください。}

\guidesec{模範解答}

\begin{answer}{演習1}
評価の例（ある製造業の中期経営計画のAI戦略について）：
\par\noindent\begin{gtabh}{colspec={Q[l,wd=22mm] X[l] X[l]}}
  観点 & 資料に根拠がある指摘 & 推測になる指摘 \\
  経営課題 & 「熟練技術者の退職による技能継承」が課題として明記されている。 & 課題の規模（退職者数、影響額）は資料にない。 \\
  使い分け & 生成AIは外部サービス、設備の異常検知は自社開発と明記。 & 自社開発の人材が確保できているかは不明。 \\
  効果の測り方 & 「2028年度に間接業務の工数20\%削減」と数値目標あり。 & 削減した工数を何に振り向けるかは書かれていない。 \\
  リスク & AIガバナンス委員会の設置が記載。 & 委員会の権限と活動内容は不明。 \\
  差別化 & 自社の設備データの蓄積が強みとされている。 & 競合も同様のデータを持つ可能性があり、差別化になるかは資料からは判断できない。 \\
\end{gtabh}
自分の評価：課題と数値目標が明確な点は評価できる。一方で、工数削減が「何のため」か（技能継承とのつながり）が弱く、
削減した時間を技能継承の活動に充てる計画があるかを確かめたい。AIの指摘は網羅的だが「資料にないこと」を推測で批判している部分があり、
それを根拠にした評価は避ける。
\end{answer}

\begin{answer}{演習2}
\begin{itemize}
  \item (A) 外部 SaaS：経費精算の問い合わせは競争優位と無関係で、固有データも社内規程程度。市販のサービスで十分。データの制約は社内規程なので小さい。
  \item (B) API 利用で自社の仕組みに組み込む（将来的に自社開発も検討）：顧客への提案は専門商社の競争優位の中核で、30年分の取引データは固有の強み。
        ただしモデルをゼロから作る人材と予算は500人規模では難しく、外部のモデルと分析基盤を API で使い、データの整備と提案のロジックに自社の資源を集中する。取引データは機密なので、データが学習に使われない契約が条件。
  \item (C) 外部 SaaS（法務向けの契約書レビューサービス）＋人の確認：契約書のチェックは重要だが差別化の源泉ではなく、専門のサービスの方が品質が高い。契約書は機密であり、データの扱いの確認が必要。最終判断は法務担当者が行う。
\end{itemize}
\end{answer}

\begin{answer}{演習3}
\par\noindent\begin{gtab}{colspec={Q[l,wd=34mm] X[l]}}
  目的と KPI & 営業担当の提案の精度を上げ、既存顧客の購買品目数を増やす。KPI：顧客あたりの年間購買品目数（現状 平均8品目→2年で10品目）、提案の成約率。 \\
  コスト & 初期：データ整備（30年分の取引データの表記統一・顧客の名寄せ）に外部委託 800万円、分析基盤と API の構築 600万円。運用：API 利用料と基盤の保守 年300万円、データ担当1名の人件費。（想定：同規模の事例と見積の相場） \\
  期待効果 & 購買品目数が2品目増えた場合、年間売上の約3\%増（現状の品目あたり売上から試算）。測りにくい効果：若手営業の提案力の底上げ。 \\
  リスクと対策 & データ整備が想定より難航する（対策：主要顧客300社から始める）。提案が外れて顧客の信頼を損なう（対策：提案は営業担当が確認してから行う、Human-in-the-loop）。取引データの外部送信（対策：学習に使われない契約、国内のサーバー）。 \\
  段階的な進め方 & 第1段階（6か月）：主要顧客300社でPoC、成約率を測る。判断基準：成約率が現状の1.3倍以上なら第2段階へ。第2段階（1年）：全顧客に展開。 \\
  投資しない場合のリスク & 競合がすでに同様の仕組みを導入しており、提案の速さと精度で劣る。熟練営業の退職で、経験に頼った提案が続かなくなる。 \\
\end{gtab}
\end{answer}

\guidesec{出席確認クイズの解説}
講義サイトの出席確認ページ（第12回）の5問です。
% 出席確認クイズ 第12回　デジタル戦略とAI戦略
%   attendance/attendance12.html から生成（解説は本ガイド用に加筆）。

\quizq{1}{DX(デジタルトランスフォーメーション)の説明として適切なものはどれでしょうか?}%
  {ホームページを作ること}%
  {紙の書類をスキャンしてPDFにすること}%
  {パソコンを新しくすること}%
  {\correct{デジタル技術でビジネスモデルや組織を変革し、競争優位を築くこと}}%
  {正解は (d)。DX（デジタルトランスフォーメーション）は、デジタル技術を使ってビジネスモデルや組織のあり方を変革し、競争優位を築くことです。(a)(b)(c) は単なるデジタル化（デジタイゼーション）であり、DX の出発点にはなりますが DX そのものではありません。}

\quizq{2}{「プラットフォーム型ビジネス」の例として適切なものはどれでしょうか?}%
  {\correct{売り手と買い手を結びつけるオンラインマーケットプレイス}}%
  {単一店舗での対面販売}%
  {自社工場での製品製造}%
  {社内の定例会議}%
  {正解は (a)。プラットフォーム型ビジネスは、売り手と買い手のように複数の利用者グループを結びつけて価値を生みます。参加者が増えるほど価値が高まるネットワーク効果が特徴です。(b)(c)(d) は単独の企業内で完結する活動です。}

\quizq{3}{AIの「自社開発」と「外部サービス利用」の比較として適切なものはどれでしょうか?}%
  {\correct{外部サービスは早く導入できるが、差別化やデータの扱いに制約がある場合がある}}%
  {両者に違いはない}%
  {自社開発は常に安い}%
  {外部サービスは常に劣る}%
  {正解は (a)。外部サービスは導入が早く手軽ですが、他社も使えるため差別化しにくく、データの扱いにも制約がある場合があります。自社開発は差別化しやすい反面、コストと人材が必要です。(c)(d) のように一方が常に優れているわけではなく、業務の重要度と自社の強みで使い分けます。}

\quizq{4}{AI投資を評価する際に考慮すべきこととして適切なものはどれでしょうか?}%
  {流行しているかどうかだけ}%
  {競合が導入したかどうかだけ}%
  {\correct{導入・運用コストと期待効果、リスク(データ・法規制など)}}%
  {初期費用だけ}%
  {正解は (c)。AI投資は、導入・運用コストと期待効果に加えて、データの確保や法規制などのリスクも含めて総合的に評価します。(a)(b) 流行や競合の動向だけで決めると目的が曖昧になり、(d) 初期費用だけでは運用コストを見落とします。}

\quizq{5}{AIによるビジネスモデル変化の例として適切なものはどれでしょうか?}%
  {価格を固定する}%
  {全員に同じ商品を郵送する}%
  {\correct{利用データをもとに、個々の顧客に合わせた提案を行うサービス}}%
  {店舗の数を増やす}%
  {正解は (c)。AIは利用データから個々の顧客の好みを推定し、パーソナライズされた提案を可能にします。これは「全員に同じものを売る」モデルから「一人ひとりに合わせる」モデルへの変化です。(a)(b)(d) はAIによるビジネスモデルの変化ではありません。}

